\usepackage[italian]{babel}
\usepackage{newlfont}
\usepackage{color}
\textwidth=450pt\oddsidemargin=0pt

\usepackage[hidelinks]{hyperref}
\usepackage{amsfonts}
\usepackage{mathtools}
\usepackage{amssymb}
\usepackage{amsthm}
\usepackage{amsmath}
\usepackage{derivative}
\usepackage{bm}
\usepackage{tasks}
\usepackage{cancel}
\usepackage{csquotes}


% bibliografia
\usepackage[
    backend=biber,
    style=numeric-comp, % Stile standard per la fisica: raggruppa [1, 2, 3] in [1-3]
    sorting=nyt,        % Ordina per Nome, Anno, Titolo
    doi=true,
    isbn=false,
    url=false
]{biblatex}

\addbibresource{bibliography.bib} % bib linking

\emergencystretch=2em


\newtheorem{theorem}{Teorema}[section]
\newtheorem{definition}[theorem]{Definizione}
\newtheorem{proposition}[theorem]{Proposizione}
\newtheorem{lemma}[theorem]{Lemma}

\theoremstyle{definition}
\newtheorem{corollary}[theorem]{Corollario}
\newtheorem{property}[theorem]{Proprietà}
\newtheorem*{remark}{Nota}